\section{Data and Methodology}

\subsection{Data Overview}

The empirical foundation is Seoul's Transportation Card Data (TCD), which
records tap-in and tap-out transactions for every transit trip. TCD captures
subway-to-subway transfers that older smart card formats miss, yielding a more
complete picture of multimodal trip chains. We analyze a single representative
weekday, February 20, 2025 (Thursday). After removing trips with no valid OTP
alternative, the dataset comprises 1,320,030 trip chains paired with 4,936,864
chain--alternative observations (3.74 alternatives per chain on average).

\begin{table}[htbp]
\centering
\tablefont
\begin{tabular}{|l|r|r|r|}
\hline
\headrow
\textbf{Chain Type} & \textbf{Chains} & \textbf{Share} & \textbf{Avg.\ Alternatives} \\
\hline
Bus Only            & 445,554   & 33.8\% & 3.51 \\
\hline
Subway Only         & 720,692   & 54.6\% & 3.82 \\
\hline
Mixed (Bus+Subway)  & 153,784   & 11.7\% & 4.01 \\
\hline
Total               & 1,320,030 & 100\%  & 3.74 \\
\hline
\end{tabular}
\caption{Trip Chain Composition}
\end{table}

For model estimation, chains with at least two alternatives (411,754) are
split 80/20 into training (329,403) and test (82,351) sets, stratified by user
type.

\begin{table}[htbp]
\centering
\tablefont
\begin{tabular}{|l|r|r|r|r|c|}
\hline
\headrow
\textbf{User Type} & \textbf{All Chains} & \textbf{Share} & \textbf{Train} &
\textbf{Test} & \textbf{Test HR Range} \\
\hline
General  & 1,044,091 & 79.1\% & 271,843 & 67,961 & 64.8--66.4\% \\
\hline
Elderly  & 177,303   & 13.4\% & 24,132  & 6,033  & 71.6--77.8\% \\
\hline
Youth    & 49,772    & 3.8\%  & 19,725  & 4,931  & 64.9--67.0\% \\
\hline
Disabled & 38,281    & 2.9\%  & 9,113   & 2,278  & 70.3--71.5\% \\
\hline
Children & 10,583    & 0.8\%  & 4,590   & 1,148  & 64.9--69.4\% \\
\hline
\end{tabular}
\caption{User Type Distribution and Test Hit Rate Range Across Models}
\end{table}

\subsection{Multi-Level Similarity Framework}

Rather than relying on a single similarity metric---a limitation noted by
Tomhave and Khani \cite{tomhave2021} in their multi-criteria transit route
choice study---we construct a hierarchical framework with four levels and ten
metrics (Table~4). Level 1 (Exact Match) is the strictest criterion, requiring
complete identity of route sequences, mode sequences, and boarding/alighting
stops. Level 2 (Structural Similarity) relaxes this to sequence-level measures
using Longest Common Subsequence (LCS) similarity. Level 3 (Stop-based
Similarity) operates at the individual stop level via Jaccard indices. Level 4
(Temporal Similarity) captures travel-time correspondence. The composite index
is:
\begin{align*}
\mathit{sim}_{\mathrm{total}} ={}& 0.30\cdot\mathit{sim}_{\mathrm{route\_seq}}
 + 0.20\cdot\mathit{sim}_{\mathrm{mode\_seq}}
 + 0.20\cdot\mathit{sim}_{\mathrm{jaccard\_full}} \\
 &+ 0.15\cdot\mathit{sim}_{\mathrm{lcs}}
 + 0.15\cdot\mathit{sim}_{\mathrm{time}}
\end{align*}

\begin{table}[htbp]
\centering
\tablefont
\begin{tabular}{|l|l|l|c|}
\hline
\headrow
\textbf{Level} & \textbf{Metric} & \textbf{Description} & \textbf{Weight} \\
\hline
1. Exact & exact\_match & Route + mode + stop identity & --- \\
\hline
\multirow{3}{*}{2. Structural}
 & sim\_route\_seq  & Route sequence LCS similarity & 0.30 \\
 & sim\_mode\_seq   & Mode sequence LCS similarity  & 0.20 \\
 & transfer\_match  & Transfer count match (binary) & --- \\
\hline
\multirow{4}{*}{3. Stop-based}
 & sim\_jaccard\_board  & Boarding stop Jaccard index  & --- \\
 & sim\_jaccard\_alight & Alighting stop Jaccard index & --- \\
 & sim\_jaccard\_full   & All-stop Jaccard index       & 0.20 \\
 & sim\_lcs             & Stop sequence LCS similarity & 0.15 \\
\hline
4. Temporal & sim\_time  & Travel time ratio similarity & 0.15 \\
\hline
Composite   & sim\_total & Weighted sum of above        & 1.00 \\
\hline
\end{tabular}
\caption{Multi-Level Similarity Framework (10 Metrics)}
\end{table}

Level 1 (Exact Match) requires complete identity of route, mode, and stop
sequences. Level 2 (Structural Similarity) uses normalized LCS similarity:
$2\cdot\mathrm{LCS}(A,B)\,/\,(|A|+|B|)$. Level 3 (Stop-based Similarity)
applies the Jaccard index
$|S_{\mathrm{obs}} \cap S_{\mathrm{alt}}|\,/\,|S_{\mathrm{obs}} \cup S_{\mathrm{alt}}|$
to boarding, alighting, and all traversed stops, complemented by stop-sequence
LCS to preserve order. Level 4 (Temporal Similarity) is defined as
$\mathit{sim}_{\mathrm{time}} = 1 - \min(1,\,
|T_{\mathrm{obs}} - T_{\mathrm{alt}}|\,/\,\max(T_{\mathrm{obs}}, T_{\mathrm{alt}}))$.
Because TCD records intermediate stops only for bus legs, subway segments use
boarding and alighting stops alone---a limitation the Jaccard formulation
handles gracefully.

\subsection{Model Specifications}

\subsubsection{Multinomial Logit (MNL)}

The utility for alternative $j$ in individual $n$'s choice set is:
\begin{align*}
V_j ={}& \beta_{\mathrm{ride}} T_{\mathrm{ride}}
 + \beta_{\mathrm{walk}} T_{\mathrm{walk}}
 + \beta_{\mathrm{transfer}} N_{\mathrm{transfer}}
 + \beta_{\mathrm{subway}} D_{\mathrm{subway}} \\
 &+ \beta_{\mathrm{peak \cdot ride}} (D_{\mathrm{peak}} \times T_{\mathrm{ride}})
 + \beta_{\mathrm{peak \cdot xfer}} (D_{\mathrm{peak}} \times N_{\mathrm{transfer}})
\end{align*}
where $T_{\mathrm{ride}}$ is in-vehicle time (min), $T_{\mathrm{walk}}$ is
total walking time (min), $N_{\mathrm{transfer}}$ is transfer count,
$D_{\mathrm{subway}}$ indicates subway inclusion, and $D_{\mathrm{peak}}$
flags peak-hour departures (07--09h, 18--20h). Waiting time was excluded after
preliminary estimation revealed multicollinearity and a counter-theoretic
positive coefficient.

\subsubsection{Mixed Logit}

To accommodate unobserved taste heterogeneity, $T_{\mathrm{walk}}$,
$N_{\mathrm{transfer}}$, and $D_{\mathrm{subway}}$ are specified as random
coefficients, $\beta_k \sim N(\mu_k, \sigma_k)$, while $T_{\mathrm{ride}}$ and
peak interactions remain fixed. Significant $\sigma_k$ values formally reject
MNL's homogeneity assumption.

\subsubsection{Latent Class Model}

Following Greene and Hensher \cite{greene2003}, a K-class model assigns each
class its own utility parameters. Peak interactions are dropped after Mixed
Logit reveals them non-significant once heterogeneity is controlled. Class
membership is modeled as a function of concomitant variables
$Z_n = [D_{\mathrm{children}}, D_{\mathrm{youth}}, D_{\mathrm{elderly}},
D_{\mathrm{disabled}}, D_{\mathrm{peak}}]$.

\subsubsection{LightGBM Benchmark}

As a non-parametric benchmark, LightGBM \cite{ke2017} classifiers predict
which alternative was chosen. Two configurations are tested: a Core model
(same four variables as MNL) and a Full model adding $D_{\mathrm{peak}}$,
user\_type, and chain composition.
